\documentclass[11pt]{article}
\usepackage[a4paper,margin=30mm]{geometry}
\usepackage[T1]{fontenc}
\usepackage{lmodern}
\usepackage{amsmath,amssymb,amsthm,mathtools}
\usepackage[hidelinks]{hyperref}
\usepackage{microtype}

\newtheorem{theorem}{Theorem}[section]
\newtheorem{lemma}[theorem]{Lemma}
\newtheorem{proposition}[theorem]{Proposition}
\newtheorem{corollary}[theorem]{Corollary}
\theoremstyle{remark}
\newtheorem{remark}[theorem]{Remark}
\newcommand{\Z}{\mathbb Z}
\newcommand{\Q}{\mathbb Q}
\newcommand{\Fp}{\mathbb F_p}
\newcommand{\chiD}{\chi_{-4}}

\title{A formal-group proof of a conjecture of Z.-H. Sun for an Ap\'ery-like sequence}
\author{Anonymous}
\date{20 September 2026}

\begin{document}
\maketitle

\begin{abstract}
Let
\[
G_n=\sum_{k=0}^n4^k\binom{2n-2k}{n-k}^{2}\binom{2k}{k}.
\]
We prove a conjecture of Z.-H. Sun asserting a three-term congruence for
$G_{(mp^r-1)/2}$ when $p\equiv1\pmod4$.  The main observation is that the
formal logarithm
\(
\ell(v)=\sum_{n\ge0}G_nv^{2n+1}/(2n+1)
\)
becomes, after the integral strict change of parameter
$v=\sqrt{\lambda(4z)}/4$, the Eichler integral of the CM newform
$\eta(4z)^6$.  The Hecke recurrence of this form and the standard
formal-group transfer theorem then give the conjectured congruence for every
odd $m$, including $p\mid m$, and every $r\ge2$.
\end{abstract}

\noindent\textbf{Keywords.} Ap\'ery-like numbers; supercongruences; formal
groups; CM modular forms; Atkin--Swinnerton-Dyer congruences.

\noindent\textbf{2020 Mathematics Subject Classification.} 11B65, 11F33,
11G20, 11A07.

\section{Introduction}

Z.-H. Sun introduced and studied the Ap\'ery-like sequence
\begin{equation}\label{eq:Gdef}
G_n=\sum_{k=0}^n4^k\binom{2n-2k}{n-k}^{2}\binom{2k}{k}.
\end{equation}
Its generating series is hypergeometric, but the indices in Sun's
Conjecture~2.7 are multiplicative: one compares coefficients whose odd
indices differ by powers of a prime.  The bridge between these two
descriptions is a one-dimensional formal group.  In the correct integral
coordinate, its logarithm is the Eichler integral of a CM eigenform, so the
three-term Hecke recurrence becomes the desired coefficient congruence.

We prove the following statement conjectured in \cite{Sun2020}.

\begin{theorem}\label{thm:main}
If $p\equiv1\pmod4$ is prime, write $p=x^2+4y^2$.  Then, for every positive
odd integer $m$ and every $r\ge2$,
\begin{equation}\label{eq:Sun}
G_{(mp^r-1)/2}\equiv(4x^2-2p)G_{(mp^{r-1}-1)/2}
-p^2G_{(mp^{r-2}-1)/2}\pmod {p^r}.
\end{equation}
\end{theorem}

The coefficient $4x^2-2p$ is therefore not an isolated numerical pattern:
it is the $p$th Hecke eigenvalue of the weight-three CM newform
$\eta(4z)^6$.  The proof has three steps.  First, the defining convolution
gives an elementary hypergeometric generating function.  Second, the strict
integral parameter $v=\sqrt{\lambda(4z)}/4$ identifies the associated
formal logarithm with the Eichler integral of that newform.  Finally, the
formal-group transfer theorem carries its Hecke recurrence back to the
coefficients $G_n$, with exactly the modulus $p^r$.  Sections~2--4 follow
these three steps in order.

The argument also gives, for $p\equiv3\pmod4$, the baseline congruence
\[
G_{(mp^r-1)/2}\equiv p^2G_{(mp^{r-2}-1)/2}\pmod {p^r}.
\]
This is weaker than the modulus $p^{2r-1}$ in Sun's Conjecture 2.6 and is
recorded only to delimit the result proved here.

\section{The generating function and its modular parameter}

Put $F(t)=\sum_{n\ge0}G_nt^n$.  We first give the elementary generating
function identity needed later.

\begin{lemma}\label{lem:gen}
One has
\begin{equation}\label{eq:gen}
F(t)=\frac{{}_2F_1(\frac12,\frac12;1;16t)}{\sqrt{1-16t}}.
\end{equation}
Consequently, with the usual modular lambda function,
\begin{equation}\label{eq:param}
F\!\left(\frac{\lambda(\tau)}{16}\right)
=\frac{\theta_3(\tau)^4}{\theta_4(\tau)^2}.
\end{equation}
\end{lemma}

\begin{proof}
The convolution in \eqref{eq:Gdef} gives
\[
F(t)=\left(\sum_{j\ge0}\binom{2j}{j}^2t^j\right)
       \left(\sum_{k\ge0}\binom{2k}{k}(4t)^k\right).
\]
The two factors are respectively
${}_2F_1(\frac12,\frac12;1;16t)$ and $(1-16t)^{-1/2}$,
which proves \eqref{eq:gen}.  Now use the classical identities
\[
{}_2F_1\!\left(\tfrac12,\tfrac12;1;\lambda(\tau)\right)
=\theta_3(\tau)^2,
\qquad
1-\lambda(\tau)=\frac{\theta_4(\tau)^4}{\theta_3(\tau)^4}.
\]
\end{proof}

Define the odd formal series
\begin{equation}\label{eq:ell}
\ell(v)=\sum_{n\ge0}\frac{G_n}{2n+1}v^{2n+1};
\qquad \ell'(v)=F(v^2).
\end{equation}
Write $q=e^{2\pi iz}$ and set
\begin{equation}\label{eq:vq}
v(q)=\frac{\sqrt{\lambda(4z)}}4
=q\prod_{n\ge1}\left(\frac{1+q^{4n}}{1+q^{4n-2}}\right)^4.
\end{equation}
The product shows that $v(q)\in q+q^3\Z[[q^2]]$; hence it is a strict
integral local parameter and its compositional inverse also belongs to
$v+v^2\Z[[v]]$.

\begin{proposition}\label{prop:log}
If $\eta(4z)^6=\sum_{n\ge1}a_nq^n$, then
\begin{equation}\label{eq:logidentity}
\ell(v(q))=\sum_{n\ge1}\frac{a_n}{n}q^n.
\end{equation}
\end{proposition}

\begin{proof}
The standard differential identity
\[
q\frac{d}{dq}\lambda(4z)
=2\theta_3(4z)^4\lambda(4z)(1-\lambda(4z))
\]
and \eqref{eq:param} give
\[
q\frac{d}{dq}\ell(v(q))
=\ell'(v(q))\,q\frac{dv(q)}{dq}
=\frac{\theta_3(4z)^4}{\theta_4(4z)^2}\,
v(q)\theta_4(4z)^4.
\]
Since $v(q)=\theta_2(4z)^2/(4\theta_3(4z)^2)$ and
$\theta_2\theta_3\theta_4=2\eta^3$, the last expression is
$\eta(4z)^6$.  Both sides of \eqref{eq:logidentity} have zero constant
term, so termwise integration proves the assertion.
\end{proof}

\section{The CM newform}

The eta product
\[
f(z)=\eta(4z)^6=\sum_{n\ge1}a_nq^n
\]
is the normalized newform of weight $3$, level $16$, and character
$\chiD$.  It has complex multiplication by $\Q(i)$.  Equivalently it is
the theta series of the Hecke character of $\Q(i)$ of infinity type $2$
and conductor $(2)$, with normalization $a_1=1$.  It follows that, for
every odd prime $p$,
\begin{equation}\label{eq:hecke}
a_{pn}-a_pa_n+\chiD(p)p^2a_{n/p}=0,
\end{equation}
where $a_{n/p}=0$ when $p\nmid n$, and
\begin{equation}\label{eq:ap}
a_p=\begin{cases}
4x^2-2p,&p=x^2+4y^2\equiv1\pmod4,\\
0,&p\equiv3\pmod4.
\end{cases}
\end{equation}
For the split case, if $\pi=x+2iy$ is chosen primary, the Hecke value is
$\pi^2+\bar\pi^2=2(x^2-4y^2)=4x^2-2p$.

\section{Formal-group transfer}

We recall the coefficient form of the standard transfer theorem; see
Li and Long \cite[Proposition 3 and Theorem 6]{LiLong}.

\begin{theorem}[formal-group coefficient transfer]\label{thm:transfer}
Let $p$ be odd and let
\[
L(q)=\sum_{n\ge1}\frac{a_n}{n}q^n,
\qquad
\ell(v)=\sum_{n\ge1}\frac{b_n}{n}v^n
\]
be related by $\ell(v(q))=L(q)$, where
$v(q)\in q+q^2\Z_{(p)}[[q]]$.  Suppose
\begin{equation}\label{eq:functional}
L(q)-\frac{A}{p}L(q^p)+B L(q^{p^2})
\in\Z_{(p)}[[q]].
\end{equation}
Then, for all positive integers $m$ and $r\ge2$,
\begin{equation}\label{eq:transfer}
b_{mp^r}-A b_{mp^{r-1}}+p^2B b_{mp^{r-2}}\equiv0\pmod {p^r}.
\end{equation}
No coprimality condition on $m$ is required.
\end{theorem}

For clarity, \eqref{eq:functional} is the functional-equation-lemma
description of a one-dimensional formal group over $\Z_{(p)}$.  Its
Cartier parameters are $\mu_0=A$, $\mu_1=-pB$, and
$\mu_j=0$ for $j\ge2$.  A strict integral change of coordinate preserves
the formal group, and the coefficient congruence in \eqref{eq:transfer}
is precisely the resulting Cartier congruence.  This also explains why the
modulus is $p^r$ and why $p\mid m$ is allowed.

\begin{proof}[Proof of the main result]
Take $L(q)=\sum a_nq^n/n$.  Relation \eqref{eq:hecke} implies
\begin{equation}\label{eq:integrality}
L(q)-\frac{a_p}{p}L(q^p)+\chiD(p)L(q^{p^2})
\in\Z_{(p)}[[q]].
\end{equation}
Indeed, comparison at an exponent divisible by $p$ is exactly
\eqref{eq:hecke}, while the remaining coefficients are $p$-integral.
By Proposition \ref{prop:log}, Theorem \ref{thm:transfer} applies through
the strict parameter \eqref{eq:vq} with $A=a_p$ and $B=\chiD(p)$.

In \eqref{eq:ell}, the coefficient $b_N$ is zero for even $N$ and equals
$G_{(N-1)/2}$ for odd $N$.  If $p\equiv1\pmod4$, then
$\chiD(p)=1$ and \eqref{eq:ap} gives $a_p=4x^2-2p$.  Substitution in
\eqref{eq:transfer}, for odd $m$, is exactly \eqref{eq:Sun}.
\end{proof}

\begin{remark}
For $p\equiv3\pmod4$, the same proof gives
\[
G_{(mp^r-1)/2}\equiv p^2G_{(mp^{r-2}-1)/2}\pmod {p^r},
\]
because $a_p=0$ and $\chiD(p)=-1$.  The sign is positive after moving the
last term in \eqref{eq:transfer} to the other side.
\end{remark}

\section*{Computational verification}

Exact rational power-series calculations were used only as diagnostics.
The accompanying program verifies through $q^{99}$ that the parameter
\eqref{eq:vq} is integral and odd and that
$q\,d\ell(v(q))/dq=\eta(4z)^6$.  It also checks \eqref{eq:ap} for every
odd prime below $100$.  None of these finite checks replaces a step of the
proof.

\section*{Use of generative artificial intelligence}

The proof presented in this article was found by OpenAI Codex.

\section*{Declarations}

\noindent\textbf{Funding.} No funding was received for this work.

\noindent\textbf{Competing interests.} The author declares no competing
interests.

\noindent\textbf{Data availability.} The exact verification program is
available at the
\href{https://github.com/huiminZheng-collab/sun-conjecture-2-7-proof-candidate}
{public GitHub repository}.

\begin{thebibliography}{9}

\bibitem{LiLong}
W.-C. W. Li and L. Long,
Atkin and Swinnerton-Dyer congruences and noncongruence modular forms,
\emph{RIMS K\^oky\^uroku Bessatsu} \textbf{B51} (2014), 269--299;
\url{https://arxiv.org/abs/1303.6228}.

\bibitem{Sun2020}
Z.-H. Sun,
New congruences involving Ap\'ery-like numbers,
preprint, arXiv:2004.07172v2 (2020),
\url{https://arxiv.org/abs/2004.07172}.

\end{thebibliography}

\end{document}
